\documentclass[11pt,compress,t,notes=noshow, aspectratio=169, xcolor=table]{beamer}

\usepackage{../../style/lmu-lecture}
% Defines macros and environments
\input{../../style/common.tex}

\title{Applied Machine Learning}
% \author{LMU}
%\institute{\href{https://compstat-lmu.github.io/lecture_iml/}{compstat-lmu.github.io/lecture\_iml}}
\date{}

\begin{document}

\newcommand{\titlefigure}{figure/performance_vs_interpretability}
\newcommand{\learninggoals}{
\item What to watch out for in practice
\item Pros and cons of performance metrics
\item Advanced evaluation practices}

\lecturechapter{Applied Performance Evaluation}
\lecture{Applied Machine Learning}


\begin{frame}{General Things to Consider in Practice}

\begin{itemize}
\setlength\itemsep{2em}
    \item Stakeholders need to understand metrics. Stakeholders are (groups of) people that share an interest in the successful completion of the use case. 
    \item Is the interpretation important? Can the results be communicated to the stakeholders? Depending on the technical expertise of stakeholders or how much time can be invested in understanding the results, performance metrics shall be chosen and evaluated carefully.
    \item What matters for the specific usecase, e.g., outlier sensitivity, cost sensitivity?
    \item What are computational limitations (e.g., time constraints)? What computational resources are at our disposal?
    \item Are we developing a model that answers the right question? What a model optimizes does not necessarily overlap with what we are trying to achieve in a specific use case.
\end{itemize}
\end{frame}

\begin{frame}{Answering the Wrong Question}
% A. W. Kimball (1957) Errors of the Third Kind in Statistical Consulting, Journal of the American Statistical Association, 52:278, 133-142
\begin{itemize}
\setlength\itemsep{1em}
    \item An approximate answer to the right question is worth much more than the right answer to the wrong question.
    \item Pay attention to but do not get lost in technical details. You might end up developing a model that answers the wrong question or optimizes the wrong metric.
    \item \textbf{Example:} An online retail company aims to increase revenue by displaying advertisements with different products, hoping that the customer is going to purchase additional products. The company measures how often the user clicks / taps on those advertisements. A model is trained to predict click rates based on features such as the product category, price, age of the customer, etc. However, this model does not actually measure additional revenue from advertisements, as many customers will click on them by accident or not go through with the purchase.
\end{itemize}
    
\end{frame}
\begin{frame}{What Affects Performance?}
% https://vuquangnguyen2016.files.wordpress.com/2018/03/applied-predictive-modeling-max-kuhn-kjell-johnson_1518.pdf

There is a multitude of technical aspects we need to consider:
\vspace{1em}
\begin{itemize}
\setlength\itemsep{1em}
    \item \textbf{Model type:} Some models are or are not suited for certain tasks. For instance, tree-based models are bad at predicting linearly shaped data.
    \item \textbf{Non-informative features:} Including features that are not correlated with the target increases the risk of modeling noise.
    \item \textbf{Pre-processing of features:} Features can or have to be processed to improve performance, depending on the chosen model.
    \item \textbf{Measurement errors / noise in the outcome or features:} Garbage in = garbage out. Our performance depends on the quality of model inputs.
    \item \textbf{Sample size:} The more data we have at our disposal, the better our algorithm is able to learn from it.
    \item \textbf{What kind of data?:} It also matters what kind of data we have. Increasing the size of the data set is not sufficient. The data need to be diverse as well.
\end{itemize}
\end{frame}



\endlecture
\end{document}
